\scanpage{79}
\mspage{35}
Si $J$ est un ``idéal'' i.e. $i\in J$, $j<i\Longrightarrow j\in J$,
alors on a $X_J^*=X_J$, donc $X_J^*$ est fermé.

Si $J$ est une ``bande'' i.e. si $i,j\in J$, $i<j$,
alors $\forall k\in I$, $(i<k<j\Longrightarrow k\in J)$,
ou ce qui revient au $=$, si $J'=\widetilde J\setminus J$ est
un ``idéal'', où $\widetilde J$ est l'idéal engendré par $J$,
alors on a
\[
(59)\qquad X_J^*=X_{\widetilde J}^*\setminus X_{J'}^*
\]
\deleted{$=X_{\unsure{\widetilde J|J}}^*$}.
\margin{différence de deux parties fermées, c'est une partie
loc\textsuperscript{t} fermée de $X$}
dont les $X_i^J\overset{\mathrm{déf}}{=}X_i\cap X_J^*$ ($i\in J$)
forment une stratification. Posons
\[
(60)\qquad V_{j_0-j_n}^J
=V_{j_0\ldots j_n}\cap X_J^*=V_{j_0-j_n}^Z
\qquad\text{pour }Z=X_J^*
\]
pour $j_0\ldots j_n\in J$, $j_0<j_1<\ldots<j_n$,
\deleted{dans la}
ce sont alors les $V_{i_0\ldots i_n}$ associés à la stratification
de type $J$ de $X_J^*=Z$. On a donc
\[
(61)\qquad\begin{aligned}
X_J^*&\simeq\underset{i,j\in J}{\unsure{\varinjlim|\uline{\lim}}}
(V_i^J,V_{ij}^J)\\
&\simeq\underset{\substack{(i_0<\ldots<i_n)\\\in\operatorname{Drap}(J)}}{\unsure{\uline{\lim}|\varinjlim}}
(V_{i_0\ldots i_n}^J)
\end{aligned}
\]
\deleted{On s'intéresse à la situation de}
\margin{78}
